\documentclass[12pt,a4paper]{article}
\usepackage[utf8]{inputenc}
\usepackage[T1]{fontenc}
\usepackage[brazil]{babel}
\usepackage{indentfirst}
\usepackage[margin=2.5cm]{geometry}

\title{Resenha Cr\'itica: Artigo 11 ---\\
The C4 Model for Visualising Software Architecture}
\author{Gustavo Pessoa Firmino Duarte}
\date{23 de Abril de 2026}

\begin{document}
\maketitle

\section{Introdu\c{c}\~ao}

O modelo C4 (\textit{Context, Containers, Components, Code}), criado por Simon
Brown, \'e uma abordagem leve e pragm\'atica para a documenta\c{c}\~ao e visualiza\c{c}\~ao
de arquiteturas de software. Diferentemente de nota\c{c}\~oes mais pesadas como
o UML completo ou o TOGAF, o C4 prop\~oe uma hierarquia de quatro n\'iveis de
abstra\c{c}\~ao, cada um direcionado a uma audi\^encia espec\'ifica: do diagrama
de contexto de alto n\'ivel, compreens\'ivel por stakeholders n\~ao t\'ecnicos,
at\'e os diagramas de c\'odigo de baixo n\'ivel, \'uteis para desenvolvedores
trabalhando na implementa\c{c}\~ao. O artigo de Brown sistematiza essa abordagem
e argumenta que a aus\^encia de documenta\c{c}\~ao arquitetural acess\'ivel \'e
uma das principais causas de decis\~oes equivocadas e dificuldades de comunica\c{c}\~ao
em equipes de desenvolvimento.

\section{Os Quatro N\'iveis de Abstra\c{c}\~ao}

O primeiro n\'ivel, o \textbf{Diagrama de Contexto} (\textit{System Context
Diagram}), mostra o sistema sendo constru\'ido e suas rela\c{c}\~oes com usu\'arios
e outros sistemas externos. \'E intencionalmente simples e pode ser apresentado
a qualquer stakeholder, t\'ecnico ou n\~ao, para estabelecer o escopo e as
fronteiras do sistema.

O segundo n\'ivel, o \textbf{Diagrama de Cont\^eineres} (\textit{Container
Diagram}), detalha as aplica\c{c}\~oes e armazenamentos de dados que comp\~oe o
sistema --- cada cont\^einer \'e uma unidade implantat\'ivel independentemente
(como uma aplica\c{c}\~ao web, uma API, um banco de dados ou uma fila de mensagens).
Esse diagrama \'e voltado a desenvolvedores e arquitetos e responde \`a pergunta:
``Como o sistema \'e dividido em blocos implant\'aveis?''

O terceiro n\'ivel, o \textbf{Diagrama de Componentes} (\textit{Component
Diagram}), detalha os componentes internos de um cont\^einer --- as unidades
l\'ogicas principais como controladores, servi\c{c}os e reposit\'orios. O quarto
n\'ivel, o \textbf{Diagrama de C\'odigo}, \'e opcional e mapeia a implementa\c{c}\~ao
em classes e interfaces, frequentemente gerado automaticamente por ferramentas
de análise est\'atica.

Brown enfatiza que o valor do C4 est\'a na consist\^encia e na legibilidade:
cada diagrama deve ser autoexplicativo, com elementos claramente rotulados quanto
a seus tipos, tecnologias e responsabilidades, sem depend\^encia de uma leg\'enda
externa. A \textit{ubiquidade} dos diagramas --- publicados em wikis, README,
conversas de onboarding --- \'e t\~ao importante quanto sua precis\~ao.

\section{Conex\~ao com a Pr\'atica Profissional}

A leitura do modelo C4 resolveu uma tens\~ao que havia sentido sem conseguir
articular: a dificuldade de saber \textit{qual n\'ivel de detalhe} usar ao
documentar uma arquitetura. No meu MVP de e-commerce, quando precisei explicar
o sistema para o professor orientador, acabei criando um diagrama informal que
misturava o n\'ivel de cont\^eineres (Node.js + SQLite + React) com detalhes de
componentes internos (rotas, servi\c{c}os, modelos Prisma). O resultado foi um
diagrama confuso que tentava dizer tudo ao mesmo tempo para audi\^encias diferentes.

Com o modelo C4, o diagrama correto para essa apresenta\c{c}\~ao seria o
\textit{Container Diagram}: React SPA $\to$ API Node.js/Express $\to$ SQLite via
Prisma. Simples, claro, sem detalhes de implementa\c{c}\~ao. Para a equipe de
desenvolvimento (eu mesmo), um \textit{Component Diagram} detalhando rotas,
middlewares e servi\c{c}os seria o mais \'util.

Na ArcelorMittal, a aus\^encia de diagramas C4 (ou qualquer documenta\c{c}\~ao
arquitetural estruturada) tornava o onboarding de novos membros do time de suporte
extremamente dependente do conhecimento t\'acito de pessoas espec\'ificas.

\section{Conclus\~ao}

O modelo C4 de Simon Brown \'e um exemplo raro de abordagem que \'e ao mesmo tempo
simples de aprender, \'util na pr\'atica e suficientemente flex\'ivel para
sistemas de qualquer escala. Sua maior contribui\c{c}\~ao \'e resolver o problema
de audi\^encia: ao definir explicitamente para quem cada n\'ivel de diagrama \'e
dirigido, o C4 elimina a tens\~ao entre diagramas de alto n\'ivel que s\~ao
abstratos demais para desenvolvedores e diagramas detalhados que s\~ao
incompreens\'iveis para stakeholders. Para qualquer equipe que sinta que sua
arquitetura est\'a mal documentada ou \'e dif\'icil de comunicar, o C4 \'e um
ponto de partida acess\'ivel e imediatamente acion\'avel.

\end{document}
